Mechanistic studies localize language-model function in parameters, heads or neurons, but those parts do not specify how one inference unfolds. We instead treat the ordered layer-state trajectory as the empirical process. Across Qwen, Llama and Gemma, realized order was 1.60--2.24 times more chord-aligned than endpoint-preserving shuffles. Random initializations were still more aligned, whereas reordered executions retained geometric regularity but lost native outputs: architecture supplies a propagation scaffold, training organizes compatible task coordinates and native order executes functionally consequential transitions. Across Pythia training, fixed task coordinates emerged only after training began and function required compatible input, block and output components. Ordered-prefix summaries did not generally predict final boundaries beyond current state. After additive interventions failed, confidence-gated local-transition actuators crossed a declared full-vocabulary boundary in 17/87, 13/87 and 9/87 held-out prompts, with 0/203, 2/203 and 0/203 non-target changes. Ordered high-dimensional states expose an observable and controllable computational process, while geometry, forecasting, action assignment and correctness remain bounded.
